\documentclass[conference]{IEEEtran}
\IEEEoverridecommandlockouts
\usepackage{cite}
\usepackage{amsmath,amssymb,amsfonts}
\usepackage{algorithmic}
\usepackage{graphicx}
\usepackage{textcomp}
\usepackage{xcolor}
\usepackage{booktabs}

\def\BibTeX{{\rm B\kern-.05em{\sc i\kern-.025em b}\kern-.08em
    T\kern-.1667em\lower.7ex\hbox{E}\kern-.125emX}}

\begin{document}
\title{Not All Items Are Created Equal\\
\vspace{0.5cm}
\large Importance-Aware Next Basket Recommendation}

\author{
\IEEEauthorblockN{Dhruv Ramesh Joshi}
\IEEEauthorblockA{\textit{IIIT-Bangalore}}
\and
\IEEEauthorblockN{R Ricky Roger}
\IEEEauthorblockA{\textit{IIIT-Bangalore}}
\and
\IEEEauthorblockN{Siddharth Prashant Kini}
\IEEEauthorblockA{\textit{IIIT-Bangalore}}
\and
\IEEEauthorblockN{Subhash Hari}
\IEEEauthorblockA{\textit{IIIT-Bangalore}}
}

\maketitle

%-----------------------------------------------------------------------
\begin{abstract}
Next basket recommendation predicts the full set of items a user will
purchase next given their history. Existing methods encode each basket
as a flat vector and treat every item as an equally important signal of
user intent. We argue this is structurally wrong: every basket is
organized around a small number of intent-defining items, with the
remainder serving as peripheral fill. We propose an architecture that
makes this distinction explicit at every stage. A BERT-style intra-basket
encoder produces per-item contextualized representations, from which a
learned importance head derives per-item importance weights initialized
from a geometric leave-one-out attribution score corrected for
population frequency. An elementwise-gated dual-stream representation
blends a full-basket summary with an importance-weighted intent centroid.
A causal GPT with rotary position embeddings models the temporal sequence
of basket representations and produces a single predicted next-basket
vector. Decoding proceeds in two conditioned stages: core items are
predicted first via a low-rank orthogonal intent projection, and
peripheral items are then predicted from a fill head whose query is
explicitly conditioned on the discovered core, bridged by a
differentiable soft intent context during training. This prevents intent
items from crowding peripheral slots while ensuring fill predictions are
coherent with the predicted intent.
\end{abstract}

\begin{IEEEkeywords}
next basket recommendation, purchase intent, importance scoring,
transformer, hierarchical decoding, sequential recommendation
\end{IEEEkeywords}


%-----------------------------------------------------------------------
\section{Introduction}
Next Basket Recommendation (NBR) is a well-established problem in the recommender systems literature, with sustained research interest over the past decade. Given a sequence of historical baskets — sets of items purchased or consumed together — the goal of an NBR system is to predict the set of items a user will select in their next interaction. This setting is prominent in real-world applications such as e-commerce and grocery shopping, where users routinely purchase multiple items in a single session.
NBR is fundamentally distinct from sequential item recommendation and session-based recommendation: unlike those settings, NBR must deal with an unordered set of multiple items as both input and output. This structural difference means that models designed for item-level recommendation are generally not directly applicable, and a dedicated body of literature has emerged to address the NBR task specifically.

A user who buys chicken, basmati rice, saffron, yogurt, and cooking oil
in one trip is not selecting five independent items. They are executing
a meal plan. The basket is the physical realization of a latent intent.
Remove the chicken and the basket loses its organizing purpose. Remove
the cooking oil and the basket remains a coherent meal.

This asymmetry is invisible to every existing next basket recommendation
(NBR) system. Standard methods encode each basket as a mean or max pool
over item embeddings and treat all items as equal contributors to the
basket's representation. The model receives identical gradient signal
whether it mispredicts the meal anchor or an incidental staple. At
inference, items are decoded by scoring all catalog items against a
single flat predicted vector and taking the top-$K$, with no mechanism
to prevent predicting four redundant beverages or to ensure that fill
items are coherent with whatever main intent was predicted.

We propose an architecture grounded in the observation that a basket is
not a flat bag but a hierarchical satisfaction of a purchase intent:
a few load-bearing items define why the basket exists, and a larger
set of fill items complete it. Our system makes this structure explicit
in the encoder, the basket representation, and the decoder.

Concretely, we contribute:
\begin{itemize}
\item A \textbf{geometric importance score} per item, measuring how much
  the basket representation shifts when that item is removed, corrected
  for population frequency to separate structural importance from rarity.
\item A \textbf{parallel importance head} on top of the BERT encoder
  that produces per-item importance weights, initialized from the
  geometric score and refined during training.
\item A \textbf{dual-stream gated basket representation} that blends the
  full-basket CLS summary with an importance-weighted intent centroid.
\item A \textbf{conditioned two-stage decoder}: core items are predicted
  from a low-rank intent projection of the GPT output; fill items are
  then predicted from a fill head whose query is conditioned on the
  discovered core via a differentiable soft bridge, preventing intent
  items from crowding fill slots while keeping the two stages jointly
  trainable.
\end{itemize}


%-----------------------------------------------------------------------
\section{Literature Survey}
\label{sec:related}

\subsection{Frequency and Neighbor Methods}

The simplest NBR baselines rely on item frequency.
P-TopFreq and G-TopFreq recommend the most personally or globally
frequent items respectively. Li et al.\cite{li2023next} propose
GP-TopFreq as a strong hybrid baseline and show that deep learning
methods frequently fail to surpass it once repeat and explore recall are
evaluated separately. TIFUKNN~\cite{hu2020modeling} and
UP-CF@r~\cite{faggioli2020recency} add temporal frequency dynamics and
recency weighting to collaborative filtering, achieving competitive
performance primarily through accurate repeat prediction.

\subsection{RNN and Graph Methods}

DREAM~\cite{yu2016dynamic} introduced RNN-based temporal modeling of
basket sequences, representing each basket as a mean pool over item
embeddings. HRM~\cite{wang2015learning} formalized the two-level
item-within-basket and basket-within-sequence hierarchy, the earliest
architectural precursor to our design. DNNTSP~\cite{yu2020predicting}
built a dynamic co-occurrence graph over items, runs graph convolution
for intra-basket item relationships, and uses masked self-attention over
the basket sequence. Despite its sophistication it assigns uniform
importance to all graph nodes. Sets2Sets~\cite{hu2019sets2sets} encodes
frequency as an explicit input feature, making it repeat-biased by
construction. Beacon~\cite{le2019correlation} captures pairwise item
correlations for coherent predictions but cannot represent directed
importance structure.

\subsection{Contrastive Methods}

CLEA~\cite{qin2021world} is the closest prior work to our importance
idea. It learns a binary relevance mask via Gumbel-Softmax to separate
relevant from irrelevant basket items and trains with a contrastive
objective. The mask is binary and learned from the prediction loss with
no geometric grounding. Our importance score is a continuous geometric
measurement computed independently of any task objective.
CL4NBR~\cite{zhu2025contrastive} adds contrastive augmentation over basket
sequences with uniform item treatment.

\subsection{Transformer Methods}

SASRec~\cite{kang2018self} demonstrated that causal self-attention
substantially outperforms RNNs for sequential recommendation.
BERT4Rec~\cite{sun2019bert4rec} applied bidirectional masked attention
with a Cloze objective to user histories and showed that bidirectional
context produces richer item representations. Our Layer 1 and Layer 2
encoders are directly inspired by these two papers respectively.
NEST is the closest architectural ancestor:
a set-within-encounter bidirectional encoder with no positional
encoding, followed by a causal sequence encoder with rotary position
embeddings. We extend this topology with importance scoring,
dual-stream fusion, and conditioned two-stage decoding.

\subsection{The Repeat-Explore Problem}

ReCANet~\cite{ariannezhad2022recanet} empirically established that
over 54\% of NBR recall on grocery datasets comes from repeat items
constituting less than 1\% of the catalog. Li et al.~\cite{li2023next}
provided the definitive evaluation, showing no published NBR method
properly balances repeat and explore prediction. They introduce separate
Repeat Recall@$K$ and Explore Recall@$K$ metrics and GP-TopFreq as the
baseline every method must beat. Our IDF-corrected importance score
directly addresses this: high-frequency items receive a population
penalty, preventing the model from treating habitual repeats as
intent-defining.

% \subsection{Attention as Importance}

% Jain and Wallace~\cite{jain2019attention} showed that learned attention
% weights are frequently uncorrelated with gradient-based importance
% measures. We do not use CLS attention weights as importance proxies.
% Instead, our geometric leave-one-out score measures the actual shift in
% the CLS output when an item is removed, grounding importance in
% the model's output behavior.


%-----------------------------------------------------------------------
\section{Motivation}
\label{sec:motivation}

\subsection{Baskets Are Organized Around Intent}

Every basket has a reason it exists. The items in a basket are not
independent purchases but a joint realization of an underlying intent:
make biryani, restock household staples, grab ingredients for a quick
weeknight meal. The intent is the ground truth. The basket is a noisy
observation of it, shaped by availability, habit, and opportunity.

This reframing has a concrete consequence. The right prediction target
is the intent underlying the next basket, from which the basket follows.
A system that predicts the basket directly, treating it as ground truth,
is predicting the noise alongside the signal with no mechanism to
distinguish them.

\subsection{Items Have Heterogeneous Importance}

Items in a basket exist on a spectrum of intent-bearing importance.
Some items define why the basket exists---removing them collapses its
purpose. Others are expected given the intent but not unique to it.
Others are peripheral and could be replaced freely without changing the
basket's meaning. A model that assigns equal gradient signal to all
three types will learn the signal and the noise simultaneously.

\subsection{Flat Decoding Fails in Two Ways}

Predicting the next basket from a single flat vector fails
structurally. First, if intent and fill scores are combined into a
single ranking, fill items may score weakly and intent items fill all
top-$K$ slots---the predicted basket contains only anchors and no
complements. Second, there is no notion of slot saturation: once a
beverage intent is satisfied, additional beverages should score low,
but a flat ranking has no mechanism for this.

The correct structure is sequential: predict intent items first, then
predict fill items conditioned on what the intent items reveal about the
basket. The fill prediction must know what was predicted as intent
before it runs.

\subsection{Frequency and Importance Are Not the Same}

Item frequency is not a reliable proxy for intent-bearing importance.
High-frequency items are not always peripheral: chicken is purchased
frequently but is genuinely load-bearing in every basket where it
appears as a main protein. Low-frequency items are not always central:
a rare imported product may be rare simply due to limited availability.
A principled importance score must capture structural contribution
independently of frequency, which motivates the IDF correction in our
geometric score.


%-----------------------------------------------------------------------
\section{Core Idea}
\label{sec:idea}

The central proposal has four parts.

\textbf{Geometric importance.} After a warmup phase, we measure each
item's structural contribution to the basket representation by computing
the Euclidean shift in the CLS output when that item is removed. This
leave-one-out displacement, averaged across all basket appearances and
corrected for population frequency via an IDF term, gives a stable
global importance score $\alpha_i^{IDF}$ per item.

\textbf{Importance head.} A lightweight MLP attached to the BERT
encoder produces per-item importance weights $\alpha_i$ from the
contextualized item representations $h_i$. It is initialized to
reproduce $\alpha_i^{IDF}$ and then trained freely, allowing it to
learn basket-specific context that the global score cannot capture.

\textbf{Dual-stream basket representation.} Two representations are
computed from each basket---a full-basket CLS summary and an
importance-weighted centroid over item representations---and fused via
a learned elementwise gate. The gate adapts as the importance head
matures: early in training it relies on the CLS summary; as $\alpha$
stabilizes it shifts toward the importance-filtered centroid.

\textbf{Conditioned two-stage decoding.} The GPT outputs a single
predicted next-basket vector. A learned low-rank orthogonal projection
extracts an intent component from this vector. Core items are predicted
from the intent component. The fill head's query is then conditioned on
the soft weighted average of the intent scores over item embeddings---a
differentiable bridge that carries the intent signal into fill
prediction without requiring a hard decoding step at training time.


%-----------------------------------------------------------------------
\section{Methodology}
\label{sec:method}

\subsection{Item Embedding Initialization}

We initialize the shared item embedding matrix
$E \in \mathbb{R}^{|V| \times d}$ via skip-gram word2vec on the basket
corpus, treating each basket as a sentence and each item as a token.
Co-occurring items are initialized nearby in $\mathbb{R}^d$. $E$ is
fully trainable throughout all phases.

\subsection{Intra-Basket Encoder with Importance Head}

Each basket $B_t = \{i_1, \ldots, i_K\}$ is encoded by a BERT-style
transformer. A learned CLS token is prepended and the sequence is
embedded:
\begin{equation}
X = [e_{\text{CLS}},\ E[i_1],\ \ldots,\ E[i_K]] \in \mathbb{R}^{(K+1) \times d}
\end{equation}

We apply $L_1$ layers of full bidirectional self-attention with no
positional encoding and no causal mask.
% \begin{align}
% A &= \operatorname{softmax}\!\left(\frac{QK^\top}{\sqrt{d}}\right),\quad
% X \leftarrow \operatorname{LN}(X + AV),\quad
% X \leftarrow \operatorname{LN}(X + \operatorname{FFN}(X))
% \end{align}

No positional encoding is used because baskets are unordered sets.
No causal mask is used because every item should attend to every other
item; intra-basket structure is fully bidirectional. The output is
$[h_{\text{CLS}}, h_1, \ldots, h_K]$.

A parallel importance head, a two-layer MLP with sigmoid output,
maps each item's contextualized representation to an importance weight:
\begin{equation}
\alpha_i = \sigma\!\left(W_2\, \operatorname{GELU}(W_1 h_i + b_1) + b_2\right) \in [0,1]
\end{equation}

The MLP operates on $h_i$ rather than on the raw embedding $E[i]$
because $h_i$ encodes item $i$ in the context of the current basket:
the same item may be more or less load-bearing depending on what else
is present. A sigmoid (not softmax) output allows independent per-item
scores rather than a zero-sum competition.

An auxiliary masked language modeling loss trains the encoder to
reconstruct randomly masked items from context, providing a direct
training signal for the intra-basket encoder independent of the
downstream temporal prediction task:
\begin{equation}
\mathcal{L}_{\text{MLM}} = \operatorname{CE}\!\left(E\, h_k^{\text{masked}},\ i_k^{\text{true}}\right)
\end{equation}

\subsection{Importance Score Initialization}

After a warmup phase, we freeze the BERT encoder and compute a
per-item global importance score. For each item $i$ in each training
basket $B_t$ we measure the shift in the CLS output when $i$ is
removed:
\begin{equation}
\delta(i, B_t) = \|h_{\text{CLS}}(B_t) - h_{\text{CLS}}(B_t \setminus \{i\})\|_2
\end{equation}

We normalize within each basket and aggregate across all appearances:
\begin{align}
\hat{\delta}(i, B_t) &= \frac{\delta(i, B_t)}{\sum_{j \in B_t} \delta(j, B_t)} \\
\bar{\delta}_i        &= \frac{1}{|\{(u,t):i \in B_t\}|} \sum_{(u,t):\,i \in B_t} \hat{\delta}(i, B_t)
\end{align}

An IDF correction separates structural importance from population
frequency:
\begin{equation}
\alpha_i^{IDF} = \bar{\delta}_i \cdot \log\frac{N}{df_i}
\end{equation}

where $N$ is the total number of training baskets and $df_i$ is the
number of baskets containing item $i$. The score is normalized to $[0,1]$.
% Items
% with fewer than $n_{\min}$ appearances have their score imputed from
% nearest neighbors in $E$.

The importance head is then pre-trained to reproduce these scores
via a mean-squared initialization loss before main training begins.

\subsection{Dual-Stream Gated Basket Representation}

Using the BERT outputs and the importance weights we construct two
basket-level representations:
\begin{align}
b_t^{\text{full}} &= h_{\text{CLS}} \\
b_t^{\text{core}} &= \frac{\sum_k \alpha_{i_k}\, h_k}{\sum_k \alpha_{i_k}}
\end{align}

$b_t^{\text{full}}$ encodes the full basket including peripheral items.
$b_t^{\text{core}}$ is the importance-weighted centroid: it emphasizes
the semantic neighborhood of load-bearing items and suppresses peripheral
ones. These two are complementary---the former is richer but noisier,
the latter is focused but incomplete.

An elementwise gate fuses them:
\begin{align}
g_t &= \sigma\!\left(W_g [b_t^{\text{full}};\, b_t^{\text{core}}]\right) \in [0,1]^d \\
b_t &= g_t \odot b_t^{\text{full}} + (1 - g_t) \odot b_t^{\text{core}}
\end{align}

The gate is elementwise rather than scalar because different dimensions
of the embedding space may be differentially reliable between the two
streams. Early in training, when $\alpha$ is poorly calibrated, the
gate relies on $b_t^{\text{full}}$; as training progresses and
importance weights stabilize, the gate shifts toward $b_t^{\text{core}}$.

\subsection{Inter-Basket Causal GPT}

The sequence $(b_1, b_2, \ldots, b_T)$ is fed into a causal GPT with
$L_2$ transformer layers. Rotary position embeddings (RoPE) are applied
to query and key vectors only, preserving each $b_t$ as a pure semantic
summary of its contents:
\begin{equation}
A = \operatorname{softmax}\!\left(\frac{\operatorname{RoPE}(Q)\,\operatorname{RoPE}(K)^\top}{\sqrt{d}} + M_{\text{causal}}\right)
\end{equation}

RoPE is used rather than learned positional embeddings because basket
sequences vary in length across users, and RoPE generalizes to
unseen lengths by encoding relative position in the attention inner
product. Loss is computed at every position $t = 1, \ldots, T-1$,
providing a training signal from the full history. The output at the
final position is $\hat{h}_{T+1} \in \mathbb{R}^d$, the predicted
next-basket representation.

\subsection{Orthogonal Intent-Fill Decomposition}

$\hat{h}_{T+1}$ contains mixed signal: some directions encode the
dominant purchase intent and others encode peripheral fill structure.
We separate them via a learned low-rank orthogonal projection
$P \in \mathbb{R}^{d \times d_k}$ ($d_k < d$) with orthonormal columns
($P^\top P = I_{d_k}$):
\begin{align}
\hat{h}^{\text{intent}} &= PP^\top \hat{h}_{T+1} \\
\hat{h}^{\text{fill}}   &= \hat{h}_{T+1} - PP^\top \hat{h}_{T+1}
\end{align}

By construction $\hat{h}^{\text{intent}} + \hat{h}^{\text{fill}} =
\hat{h}_{T+1}$ exactly, and the two components are orthogonal.
Orthonormality is enforced via periodic Gram-Schmidt
re-orthonormalization. An orthogonality regularizer
$\mathcal{L}_{\text{orth}} = \|\hat{h}^{\text{intent}} \cdot
\hat{h}^{\text{fill}}\|^2$ prevents drift toward the identity during
training.

\subsection{Conditioned Two-Stage Decoding}
\label{subsec:decode}

The core design choice in our decoder is that fill prediction is
conditioned on intent prediction, not run in parallel from a combined
score. Combining intent and fill scores additively would allow intent
items to dominate all $K$ output slots when fill scores are weak.
Instead we predict $K_1$ core items first, then predict $K_2$ fill
items whose query has been shifted by the discovered intent.

\paragraph{Stage 1 --- Core item prediction.}
Items are scored against the intent component:
\begin{equation}
s_i^{\text{intent}} = e_i \cdot \hat{h}^{\text{intent}}
\end{equation}

During training we construct a soft intent context vector using a
temperature-scaled softmax over all item scores---a differentiable
proxy for the hard top-$K_1$ selection:
\begin{equation}
\tilde{c} = \sum_{i \in V} \operatorname{softmax}_\tau(s^{\text{intent}})_i \cdot e_i \in \mathbb{R}^d
\end{equation}
 At inference
we run a residual loop: score all items against $\hat{h}^{\text{intent}}$,
pick the best, subtract its projected contribution, repeat $K_1$ times:
\begin{align}
\hat{i}_k &= \operatorname*{arg\,max}_{i \notin \hat{C}}\; e_i \cdot r_{k-1}^{\text{intent}} \\
r_k^{\text{intent}} &= r_{k-1}^{\text{intent}} - PP^\top e_{\hat{i}_k}
\end{align}

The residual loop enforces natural diversity: once an item is picked and
subtracted, redundant items score low against the updated residual.

\paragraph{Stage 2 --- Fill item prediction conditioned on intent.}
The fill query is the fill component shifted by the soft intent context:
\begin{equation}
\hat{h}^{\text{fill} \mid \text{intent}} = \operatorname{LN}\!\left(\hat{h}^{\text{fill}} + W_{\text{cond}}\, \tilde{c}\right)
\end{equation}

$W_{\text{cond}} \in \mathbb{R}^{d \times d}$ is learned end-to-end.
This shifts the fill query toward the embedding neighborhood of
whatever intent was discovered, so fill items are predicted as
complements to the specific discovered intent rather than in the
abstract. At inference we use $\tilde{c} = \bar{e}_{\hat{C}} =
\frac{1}{K_1}\sum_{k} e_{\hat{i}_k}$ (mean of predicted core
embeddings) and run the residual loop in the fill subspace:
\begin{align}
\hat{i}_{K_1+k} &= \operatorname*{arg\,max}_{i \notin \hat{C}}\; e_i \cdot r_{k-1}^{\text{fill}} \\
r_k^{\text{fill}} &= r_{k-1}^{\text{fill}} - (I - PP^\top) e_{\hat{i}_{K_1+k}}
\end{align}

The differentiable bridge $\tilde{c}$ means that gradients from
$\mathcal{L}_{\text{fill}}$ flow back through $W_{\text{cond}}$,
through $\tilde{c}$, and into $s^{\text{intent}}$, jointly training
both stages. The intent head is incentivized to produce representations
that are useful not only for core prediction but also for conditioning
fill prediction.

\subsection{Training Objective}

The ground-truth basket is partitioned using a threshold $\tau_\alpha$
on $\alpha^{IDF}$:
\begin{equation}
C^*_{T+1} = \{i \in B_{T+1} : \alpha_i^{IDF} > \tau_\alpha\},\quad
F^*_{T+1} = B_{T+1} \setminus C^*_{T+1}
\end{equation}

The intent loss is weighted by $\alpha^{IDF}$ so that missing
load-bearing items incurs a larger penalty:
\begin{equation}
\mathcal{L}_{\text{intent}} = -\frac{1}{T}\sum_t \sum_j \alpha_j^{IDF}
  \bigl[y_{t,j}\log\sigma(s_j^{\text{intent}})
       + (1-y_{t,j})\log(1-\sigma(s_j^{\text{intent}}))\bigr]
\end{equation}

The fill loss is computed on the conditioned fill scores with uniform
item weights:
\begin{equation}
\mathcal{L}_{\text{fill}} = -\frac{1}{T}\sum_t
  \sum_{j \in F^*_{t+1}}
  \bigl[y_{t,j}\log\sigma(s_j^{\text{fill}\mid\text{intent}})
       + (1-y_{t,j})\log(1-\sigma(s_j^{\text{fill}\mid\text{intent}}))\bigr]
\end{equation}

The total loss is:
\begin{equation}
\mathcal{L} = \mathcal{L}_{\text{intent}}
            + \lambda\,\mathcal{L}_{\text{fill}}
            + \gamma\,\mathcal{L}_{\text{orth}}
            + \eta\,\mathcal{L}_{\text{MLM}}
\end{equation}

\subsection{Training Schedule}

Training proceeds in four phases. In \textbf{Phase 0}, $E$ is
initialized via skip-gram. In \textbf{Phase 1} (warmup), BERT, GPT,
and the gated fusion are trained with uniform BCE; the importance head
is frozen. In \textbf{Phase 2}, BERT is frozen and $\alpha^{IDF}$
scores are computed via the leave-one-out passes; the importance head
is pre-trained to reproduce them. In \textbf{Phase 3}, all components
are unfrozen and trained jointly with the full loss.


%-----------------------------------------------------------------------
\section{Inference}
\label{sec:inference}

At inference time no gradients are computed and the soft intent context
$\tilde{c}$ is replaced by its hard counterpart. The full procedure
for a user with basket history $(B_1, \ldots, B_T)$ is as follows.

\paragraph{Step 1 --- Encode history.}
Each past basket $B_j$ is passed through the BERT encoder and
importance head to produce $h_{\text{CLS}}^j$, $\{h_k^j\}$, and
$\{\alpha_{i_k}^j\}$. The dual-stream gated fusion yields one
basket vector per timestep:
\begin{equation}
b_j = g_j \odot h_{\text{CLS}}^j + (1 - g_j) \odot b_j^{\text{core}},
\quad j = 1, \ldots, T
\end{equation}

\paragraph{Step 2 --- Predict next basket representation.}
The sequence $(b_1, \ldots, b_T)$ is passed through the causal GPT.
The output at the final position is the predicted next-basket
representation:
\begin{equation}
\hat{h}_{T+1} = \text{GPT}(b_1, \ldots, b_T) \in \mathbb{R}^d
\end{equation}

\paragraph{Step 3 --- Decompose into intent and fill.}
The learned projection $P$ splits $\hat{h}_{T+1}$ into orthogonal
components:
\begin{equation}
\hat{h}^{\text{intent}} = PP^\top \hat{h}_{T+1}, \qquad
\hat{h}^{\text{fill}}   = \hat{h}_{T+1} - PP^\top \hat{h}_{T+1}
\end{equation}

\paragraph{Step 4 --- Core item decoding via residual loop.}
Initialize $r_0^{\text{intent}} = \hat{h}^{\text{intent}}$ and
$\hat{C} = \varnothing$. For $k = 1, \ldots, K_1$:
\begin{align}
\hat{i}_k &= \operatorname*{arg\,max}_{i \notin \hat{C}}\;
             e_i \cdot r_{k-1}^{\text{intent}} \\
r_k^{\text{intent}} &= r_{k-1}^{\text{intent}} - PP^\top e_{\hat{i}_k} \\
\hat{C} &\leftarrow \hat{C} \cup \{\hat{i}_k\}
\end{align}

Each selected item's projected embedding is subtracted from the
residual before the next query. Items that point in the same direction
as an already-selected item score low against the updated residual,
preventing redundant picks without any explicit diversity constraint.

\paragraph{Step 5 --- Condition fill query on discovered core.}
Compute the mean embedding of the predicted core set:
\begin{equation}
\bar{e}_{\hat{C}} = \frac{1}{K_1} \sum_{k=1}^{K_1} e_{\hat{i}_k}
\end{equation}

Shift the fill component toward the discovered intent neighborhood:
\begin{equation}
\hat{h}^{\text{fill} \mid \text{intent}}
  = \operatorname{LN}\!\left(\hat{h}^{\text{fill}}
    + W_{\text{cond}}\,\bar{e}_{\hat{C}}\right)
\end{equation}

This is the hard inference counterpart of the soft bridge $\tilde{c}$
used during training. The fill query now encodes not just the generic
fill signal from $\hat{h}^{\text{fill}}$ but also the specific
embedding context of what the intent stage actually predicted.

\paragraph{Step 6 --- Fill item decoding via conditioned residual loop.}
Initialize $r_0^{\text{fill}} = \hat{h}^{\text{fill} \mid \text{intent}}$.
For $k = 1, \ldots, K_2$:
\begin{align}
\hat{i}_{K_1+k} &= \operatorname*{arg\,max}_{i \notin \hat{C}}\;
                   e_i \cdot r_{k-1}^{\text{fill}} \\
r_k^{\text{fill}} &= r_{k-1}^{\text{fill}}
                   - (I - PP^\top) e_{\hat{i}_{K_1+k}} \\
\hat{C} &\leftarrow \hat{C} \cup \{\hat{i}_{K_1+k}\}
\end{align}

Fill item embeddings are projected into the fill subspace before
subtraction to keep the residual within that subspace. Items already
selected in Stage 1 are excluded from the candidate set.

\paragraph{Final output.}
The predicted basket is:
\begin{equation}
\hat{B}_{T+1} = \hat{C}, \quad |\hat{B}_{T+1}| = K_1 + K_2
\end{equation}

The two hyperparameters $K_1$ and $K_2$ control the size of the core
and fill sets respectively. In practice $K_1$ is small (typically
1--3) and $K_2$ is set to match the target basket size minus $K_1$.
Both can be tuned on a validation set or set per-user based on
historical average basket size.

The entire inference procedure reduces to matrix--vector multiplications
and nearest-neighbor lookups in $E$, which can be accelerated with
approximate nearest-neighbor indices (e.g., FAISS) over the full
catalog. The computational cost is $O((K_1 + K_2) \cdot |V|)$ scoring
operations after the BERT and GPT forward passes, linear in catalog
size and basket size.


%-----------------------------------------------------------------------
\section{Conclusion}
\label{sec:conclusion}

We presented an intent-aware next basket recommendation system that
makes explicit the asymmetry between load-bearing and peripheral items
present in every basket. A geometric importance score, corrected for
population frequency, initializes a learned importance head that shapes
both the basket representation and the training signal. A dual-stream
gated encoder produces a basket representation that blends full-basket
context with an importance-filtered intent centroid. A causal GPT models
basket sequences and produces a predicted next-basket vector, which a
low-rank orthogonal projection decomposes into intent and fill
components. Decoding proceeds in two conditioned stages connected by a
differentiable soft intent context, preventing intent items from crowding
fill slots while keeping the two stages jointly trained.

Future work will evaluate on standard benchmarks (Dunnhumby, Instacart,
TaFeng) with separate Repeat Recall@$K$ and Explore Recall@$K$ metrics
and will explore extending the two-level decomposition to finer
hierarchical intent structures.

%-----------------------------------------------------------------------
\bibliographystyle{IEEEtran}
\bibliography{refs}

\end{document}
